% ============================================================
%  Softinsa Badge Platform — Relatório
%  Template com as cores do projeto
% ============================================================
\documentclass[12pt,a4paper]{report}

% --- Encoding & Language ---
\usepackage[utf8]{inputenc}
\usepackage[T1]{fontenc}
\usepackage[portuguese]{babel}

% --- Geometry ---
\usepackage[
  top=2.5cm,
  bottom=2.5cm,
  left=3cm,
  right=2.5cm
]{geometry}

% --- Fonts ---
\usepackage{mathptmx}
\usepackage{microtype}
\usepackage{indentfirst}

% --- Line Spacing ---
\usepackage{setspace}
\onehalfspacing

% --- Colors (from project palette) ---
\usepackage{xcolor}
\definecolor{primary}{HTML}{00B8E0}        % --color-primary
\definecolor{secondary}{HTML}{39639C}      % --color-secondary
\definecolor{onprimary}{HTML}{FFFFFF}      % --color-on-primary
\definecolor{primarycontainer}{HTML}{B9EBF6}% --color-primary-container
\definecolor{background}{HTML}{FDF7FF}     % --color-background
\definecolor{onsurface}{HTML}{171C1E}      % --color-on-surface
\definecolor{outline}{HTML}{70787C}        % --color-outline

% --- Graphics & Floats ---
\usepackage{graphicx}
\usepackage{float}
\usepackage{caption}
\captionsetup{
  font=small,
  labelfont={bf,color=secondary},
  textfont={color=onsurface}
}

% --- Tables ---
\usepackage{booktabs}
\usepackage{tabularx}

% --- Lists ---
\usepackage{enumitem}

% --- Hyperlinks ---
\usepackage{hyperref}
\hypersetup{
  colorlinks=true,
  linkcolor=secondary,
  citecolor=secondary,
  urlcolor=primary,
  bookmarksnumbered=true
}

% --- Headers & Footers ---
\usepackage{fancyhdr}
\pagestyle{fancy}
\fancyhf{}
\fancyhead[L]{\small\color{outline}\leftmark}
\fancyhead[R]{\small\color{outline}\thepage}
\renewcommand{\headrulewidth}{0.4pt}
\renewcommand{\headrule}{\hbox to\headwidth{\color{primary}\leaders\hrule height \headrulewidth\hfill}}
\fancyfoot[C]{\small\color{outline}Softinsa Badge Platform}

% --- Chapter & Section Styling ---
\usepackage{titlesec}

\titleformat{\chapter}[display]
  {\normalfont\Huge\bfseries\color{secondary}}
  {\colorbox{primary}{\parbox{1.5cm}{\centering\color{onprimary}\fontsize{48}{52}\selectfont\thechapter}}}
  {20pt}
  {}

\titleformat{\section}
  {\normalfont\Large\bfseries\color{secondary}}
  {\thesection}{1em}{}
  [\vspace{2pt}{\color{primary}\hrule height 1.5pt}]

\titleformat{\subsection}
  {\normalfont\large\bfseries\color{secondary!80!black}}
  {\thesubsection}{1em}{}

% --- Code Listings ---
\usepackage{listings}
\lstset{
  basicstyle=\ttfamily\small,
  backgroundcolor=\color{primarycontainer!30},
  frame=single,
  rulecolor=\color{primary!50},
  keywordstyle=\color{secondary}\bfseries,
  commentstyle=\color{outline}\itshape,
  stringstyle=\color{primary!70!black},
  breaklines=true,
  tabsize=2
}

% --- Bibliography ---
\usepackage[style=numeric,sorting=nyt,backend=biber]{biblatex}
% \addbibresource{referencias.bib}  % Uncomment when bib file is ready

% --- Glossary (optional) ---
% \usepackage[toc]{glossaries}
% \makeglossaries

% --- Custom Commands ---
\newcommand{\keyword}[1]{\textbf{\color{primary}#1}}
\newcommand{\role}[1]{\textit{\color{secondary}#1}}

% ============================================================
\begin{document}

% --- Cover Page ---
\begin{titlepage}
  \begin{center}

    \vspace*{0.5cm}

    % Logos
    \includegraphics[height=2.5cm]{images/logo-escola.png}
    \hspace{2cm}
    \includegraphics[height=2.5cm]{images/logo-softinsa.png}\\[1.5cm]

    {\color{primary}\rule{\textwidth}{2pt}}\\[0.6cm]

    {\fontsize{26}{32}\selectfont\bfseries\color{secondary}
      Plataforma de Badges da Softinsa}\\[0.5cm]

    {\fontsize{14}{18}\selectfont\color{onsurface}
      Relatório da Unidade Curricular de Projeto Integrado}\\[0.3cm]

    {\color{primary}\rule{\textwidth}{2pt}}\\[1.2cm]

    {\large\color{onsurface}
      Curso de Licenciatura em Engenharia Informática\\[0.3cm]
      Ano Letivo: 2.º}\\[1.8cm]

    {\Large\color{secondary}\bfseries
      Hugo Alexandre Pereira Afonso} {\large\color{outline}(pv30032)}\\[0.3cm]
    {\Large\color{secondary}\bfseries
      Mateus Valente Frias Silva} {\large\color{outline}(pv29989)}\\[0.3cm]
    {\Large\color{secondary}\bfseries
      Rodrigo de Almeida Martins} {\large\color{outline}(pv30773)}\\[0.3cm]
    {\Large\color{secondary}\bfseries
      Susana Daniela Tavares} {\large\color{outline}(pv27467)}\\[2cm]

    {\large\color{onsurface}
      Trabalho efetuado sob a orientação de:\\[0.3cm]
      \textbf{Artur Sousa}\\[0.15cm]
      \textbf{Júlio Florentino}}

    \vfill

    {\large\color{outline}
      Junho 2026}

  \end{center}
\end{titlepage}

% --- Numbering ---
\pagenumbering{roman}

% ============================================================
% Agradecimentos
% ============================================================
\chapter*{Agradecimentos}
\addcontentsline{toc}{chapter}{Agradecimentos}

Enquanto grupo, gostaríamos de manifestar o nosso sincero agradecimento aos professores Artur Sousa e Júlio Florentino pelo apoio, orientação e tempo que nos dedicaram ao longo deste trabalho.

Gostaríamos também de expressar a nossa gratidão à Softinsa, bem como a todos os que a representaram, principalmente o engenheiro Rui Costa, pelo tempo investido na elaboração do enunciado deste projeto.

\newpage

% ============================================================
% Resumo
% ============================================================
\chapter*{Resumo}
\addcontentsline{toc}{chapter}{Resumo}

% Escrever resumo aqui.

\vspace{1cm}
\noindent\textbf{Palavras-chave:} % gamificação, badges digitais, credenciais, ...

\newpage

% ============================================================
% Abstract (optional)
% ============================================================
\chapter*{Abstract}
\addcontentsline{toc}{chapter}{Abstract}

% Write abstract here.

\vspace{1cm}
\noindent\textbf{Keywords:} % gamification, digital badges, credentials, ...

\newpage

% ============================================================
% Lista de Abreviaturas
% ============================================================
\chapter*{Lista de Abreviaturas}
\addcontentsline{toc}{chapter}{Lista de Abreviaturas}

\begin{tabbing}
  \hspace{3cm}\= \kill
  \textbf{API}  \> Application Programming Interface \\[4pt]
  \textbf{AWS}  \> Amazon Web Services \\[4pt]
  \textbf{BD}   \> Base de Dados \\[4pt]
  \textbf{IBM}  \> International Business Machines \\[4pt]
  \textbf{JWT}  \> JSON Web Token \\[4pt]
  \textbf{KPI}  \> Key Performance Indicator \\[4pt]
  \textbf{LP}   \> Learning Path \\[4pt]
  \textbf{RGPD} \> Regulamento Geral sobre a Proteção de Dados \\[4pt]
  \textbf{SL}   \> Service Line \\[4pt]
  \textbf{SLA}  \> Service Level Agreement \\[4pt]
  \textbf{SLL}  \> Service Line Leader \\[4pt]
  \textbf{TM}   \> Talent Manager \\[4pt]
  \textbf{UC}   \> Unidade Curricular \\[4pt]
  \textbf{UI}   \> User Interface \\[4pt]
  \textbf{URL}  \> Uniform Resource Locator \\[4pt]
  \textbf{UX}   \> User Experience \\
\end{tabbing}

\newpage

% ============================================================
% Glossário
% ============================================================
\chapter*{Glossário}
\addcontentsline{toc}{chapter}{Glossário}

\begin{tabularx}{\textwidth}{@{}l@{\hspace{1em}}X@{}}
  \textbf{\color{secondary}Badge}                & Credencial digital verificável que atesta a obtenção de um determinado nível de competências ou a conclusão de requisitos específicos por parte do consultor. \\[6pt]
  \textbf{\color{secondary}Consultor}             & Perfil de utilizador da plataforma que representa o colaborador/técnico da Softinsa responsável por submeter evidências e progredir nos itinerários de aprendizagem. \\[6pt]
  \textbf{\color{secondary}Credly}                & Plataforma global de referência no mercado para a emissão, gestão e partilha de badges e certificações digitais. \\[6pt]
  \textbf{\color{secondary}Evidência}             & Documento digital (certificado, diploma, relatório em formato PDF, imagem ou outro) submetido pelo consultor que comprova a realização de um determinado requisito. \\[6pt]
  \textbf{\color{secondary}Gamificação}           & Aplicação de dinâmicas, mecânicas e elementos típicos de jogos (como sistema de pontos, conquistas especiais e rankings) num contexto corporativo para motivar a aprendizagem. \\[6pt]
  \textbf{\color{secondary}Learning Path}         & Estrutura macro de formação contínua dentro da empresa, que agrupa várias SLs e áreas técnicas. \\[6pt]
  \textbf{\color{secondary}Requisito}             & Meta ou curso específico associado a um determinado nível, cuja conclusão e validação são obrigatórias para obter o respetivo badge. \\[6pt]
  \textbf{\color{secondary}Service Line}          & Divisão interna da empresa que agrupa áreas de competências técnica afins. \\[6pt]
  \textbf{\color{secondary}Softinsa}              & Empresa do grupo IBM em Portugal, focada em serviços de consultoria tecnológica e gestão de aplicações, que forneceu o tema para este projeto integrado. \\[6pt]
  \textbf{\color{secondary}Workflow de Estados}   & Ciclo de vida sequencial pelo qual uma candidatura a um badge transita no sistema (compreendendo os estados `Open', `Submitted', `Em Validação' e `Fechado') consoante as ações dos avaliadores. \\
\end{tabularx}

\newpage

% ============================================================
% Table of Contents, Figures, Tables
% ============================================================
\tableofcontents
\newpage
\listoffigures
\newpage
\listoftables
\newpage

% --- Switch to arabic numbering ---
\pagenumbering{arabic}

% ============================================================
% Introdução
% ============================================================
\chapter{Introdução}
\label{ch:introducao}

O presente relatório documenta todo o processo de idealização, planeamento e implementação da Plataforma de Badges da Softinsa, a propósito da Unidade Curricular (UC) de Projeto Integrado. O tema do projeto foi concebido e fornecido pela Softinsa, com o propósito de desafiar os estudantes a desenvolverem uma solução tecnológica robusta e adaptada a uma necessidade real do contexto corporativo da empresa. Sendo o desenvolvimento deste ecossistema efetuado integralmente a partir do zero, este documento detalha as várias etapas do ciclo de vida do software, compreendendo a modelação conceptual do modelo de dados, a conceção da experiência e interface do utilizador (UI/UX) através de protótipos, e a subsequente codificação das três frentes da solução: a API REST de backend, o portal Web e a aplicação Mobile (apenas para o perfil de consultor).

\section{Enquadramento e Motivação}

No atual panorama do setor das tecnologias da informação, a competitividade das empresas está intrinsecamente ligada à capacidade de valorização e atualização contínua do seu capital humano. Como empresa líder em consultoria tecnológica, a Softinsa reconhece que a formação contínua e a certificação de competências dos seus colaboradores são pilares fundamentais para sustentar o crescimento do negócio, responder às exigências do mercado e manter equipas de alto rendimento. Para organizar este esforço de capacitação, a empresa dispõe de múltiplos Learning Paths (LP), ou itinerários de aprendizagem, estruturados por Service Lines (SL) e áreas técnicas específicas, oferecendo uma progressão de carreira clara através de vários níveis de consolidação de conhecimento.

Contudo, a gestão manual e descentralizada destas formações -- frequentemente realizadas em plataformas externas como a Udemy, IBM, AWS ou Microsoft -- acarreta desafios operacionais significativos. Identificou-se uma clara dificuldade em evidenciar e validar de forma ágil as competências efetivamente adquiridas pelos consultores. Para responder a isto, o ecossistema da Softinsa exige um fluxo rigoroso de auditoria assente em papéis bem definidos: o Consultor, que executa as formações e submete as suas evidências; o Talent Manager (TM), responsável por uma primeira validação global e transversal dos documentos submetidos; e o Service Line Leader (SLL), que efetua a validação técnica e final com base na sua área específica de especialização. Todo este ecossistema é centralizado e parametrizado pelo perfil de Administrador, que faz a gestão de conteúdos e utilizadores da plataforma.

Atualmente, a ausência de um sistema automatizado e padronizado que una estes perfis não só reduz a visibilidade interna do leque de competências disponíveis, como também limita o impacto reputacional que certificações verificáveis poderiam trazer junto de clientes e parceiros. Aliado a isto, a falta de uma dinâmica de gamificação no ecossistema corporativo traduz-se numa oportunidade perdida para motivar os colaboradores a assumirem um papel mais proativo na sua própria evolução profissional. Os consultores necessitam de estímulos visuais claros, reconhecimento de marcos alcançados e uma perceção tangível do seu progresso dentro da organização.

Foi neste contexto que surgiu a necessidade de desenvolver a presente plataforma de badges digitais. Inspirada em soluções de referência como o Credly.com, esta ferramenta visa colmatar as falhas identificadas através da centralização do fluxo de submissão e auditoria de evidências entre os diferentes intervenientes, da introdução de mecânicas de recompensa por pontos e do fortalecimento de uma cultura colaborativa e transparente dentro da Softinsa.

\section{Objetivos}

O objetivo geral deste projeto consiste no desenvolvimento, a partir do zero, de um ecossistema tecnológico integrado para a gestão de badges digitais da Softinsa, que replique e adapte as boas práticas de plataformas de referência de mercado. A solução visa aproximar a empresa da sua rede de colaboradores de forma colaborativa e ágil, racionalizando os processos de validação de competências.

Para concretizar este propósito, foram definidos os seguintes objetivos específicos (para além dos fornecidos pela própria Softinsa no enunciado), segmentados por vertentes técnicas e operacionais:

\subsection{Desenvolvimento Tecnológico e Arquitetural}

\begin{itemize}
  \item \textbf{Conceção do Modelo de Dados:} Desenhar e implementar uma base de dados estruturada e escalável, capaz de suportar o fluxo de estados das candidaturas e que permita, futuramente, a expansão para novos LPs (como o de Power Skills);
  \item \textbf{Prototipagem:} Desenvolvimento de um Protótipo de Alta Fidelidade (PAF) de forma a explorar a melhor solução visual (UI) e também no que toca à experiência do utilizador (UX);
  \item \textbf{Implementação da API REST (Backend):} Desenvolver o motor lógico central do sistema para garantir a segurança dos dados, a autenticação robusta de utilizadores e a comunicação fluida com as interfaces;
  \item \textbf{Construção do Portal Web:} Desenhar e codificar um portal responsivo dedicado aos perfis de gestão (Administrador, TM, SLL) e também de consulta para o consultor, com foco na apresentação de métricas, relatórios de auditoria e ferramentas de parametrização;
  \item \textbf{Desenvolvimento da Aplicação Mobile:} Desenvolver uma aplicação móvel focada na experiência do consultor, permitindo-lhe gerir o seu perfil, acompanhar o seu progresso visual e submeter evidências de formação de forma intuitiva.
\end{itemize}

\subsection{Funcionalidades Operacionais e de Negócio}

\begin{itemize}
  \item \textbf{Automatização do Workflow de Aprovação:} Digitalizar o ciclo de vida de uma candidatura, garantindo o registo auditável e o histórico de decisões e comentários entre o consultor, TM e SLL;
  \item \textbf{Introdução de Mecânicas de Gamificação:} Implementar um sistema motivacional assente na atribuição de pontos por badges obtidos e na celebração de conquistas especiais para incentivar a formação contínua;
  \item \textbf{Garantia de Transparência e Verificabilidade:} Criar um sistema de verificação pública através de ligações para cada badge, permitindo aos consultores partilharem as suas competências certificadas de forma credível em plataformas como o LinkedIn.
\end{itemize}

\section{Estrutura do Relatório}

% Descrever a organização dos capítulos.

% ============================================================
% (Additional chapters go here)
% ============================================================

% \chapter{Estado da Arte}
% \chapter{Análise de Requisitos}
% \chapter{Arquitetura e Desenho}
% \chapter{Implementação}
% \chapter{Testes e Validação}

% ============================================================
% Conclusão
% ============================================================
\chapter{Conclusão}
\label{ch:conclusao}

\section{Síntese do Trabalho}

% Sintetizar o trabalho realizado.

\section{Trabalho Futuro}

% Descrever possíveis melhorias e desenvolvimentos futuros.

% ============================================================
% Bibliografia
% ============================================================
% \printbibliography[heading=bibintoc,title={Referências}]

% ============================================================
% Appendices (optional)
% ============================================================
% \appendix
% \chapter{Diagramas Complementares}

\end{document}
